\documentclass[11pt,a4paper]{article}
\usepackage[UTF8]{ctex}
\usepackage{geometry}
\geometry{left=2.5cm,right=2.5cm,top=2.5cm,bottom=2.5cm}
\usepackage{amsmath,amssymb,amsthm}
\usepackage{hyperref}
\usepackage{booktabs}
\usepackage{enumitem}
\usepackage{xltabular}
\hypersetup{colorlinks=true,linkcolor=blue,urlcolor=blue}

\newtheorem{definition}{定义}[section]
\newtheorem{theorem}{定理}[section]
\newtheorem{proposition}{命题}[section]
\newtheorem{corollary}{推论}[section]
\newtheorem{remark}{注记}[section]

\title{基于生成论的网络安全分析能力调度系统\\
第2卷：对抗性安全与信息战}
\author{李龙胜}
\date{\today}

\begin{document}

\maketitle

\begin{center}
\textit{
算法可以公开，但参数必须是秘密。\\
借鉴密码学中AES的设计哲学，将安全性收敛至核心参数。\\
本卷建立防御方如何隐藏参数、反制探测的完整理论。\\
从柯克霍夫框架到$\varepsilon$-不可区分性的严格证明，\\
再到探测与反探测的博弈论模型。\\
}
\end{center}

\vspace{5mm}

\begin{abstract}
本卷的核心目标是：在算法完全公开的前提下，通过参数保密来构建防御方的安全边界。
2.1节建立三层架构（公开层、秘密层、扰动层），完成AES设计哲学在安全运营中的迁移。
2.2节借鉴密码学IND-CPA范式，建立$\varepsilon$-参数不可区分性的严格安全定义，
并给出基于信息论下界、欠定推断和跨周期累积无效性的双层保障证明。
2.3节将安全分析上升为信息战，建立探测与反探测博弈模型，
分析蓝队最优伪造策略与蜜罐式反探测战术。
\end{abstract}

\tableofcontents

\section{2.1 柯克霍夫框架：算法公开与参数保密}

\subsection{核心原则的迁移}

密码学中，一个密码系统应该做到：即便算法完全公开，只要密钥是秘密的，系统就仍然是安全的（柯克霍夫原则）。我们将这一原则迁移至安全运营领域：

\begin{center}
\textbf{算法透明化，参数密钥化。}\\
系统的公信力建立在算法的公开审查之上，\\
系统的对抗安全收敛至参数的秘密性之上。
\end{center}

迁移对照表：
\begin{center}
\begin{tabular}{c|c}
\toprule
\textbf{密码学} & \textbf{告警分级算法} \\
\midrule
加密算法（公开） & 择优公式、分析率计算、配额分配算法（公开） \\
密钥（秘密） & 核心参数：$A_i^{\min}$、$\alpha_i$、扰动幅度等（秘密） \\
密文（可被截获） & 告警流和裁定结果（攻击者可观察） \\
已知明文攻击 & 攻击者用已知攻击试探系统响应，试图反推参数 \\
\bottomrule
\end{tabular}
\end{center}

\subsection{三层架构}

系统被垂直划分为三个独立的层，以实现算法透明与参数保密的解耦。

\begin{definition}[公开层]
择优公式、配额分配规则、反馈更新类型（翻倍/衰减/下限保护）全部公开，接受学术界和安全社区的审查与攻击验证。
\end{definition}

\begin{definition}[秘密层]
以下核心参数不公开，仅存储于SOC的加密参数库中：
\begin{itemize}
    \item 每类告警的损失下限 \(A_i^{\min}\)
    \item 敏感度系数 \(\alpha_i\) 的精确值
    \item 抗投毒衰减中使用的外部威胁情报源权重
    \item 翻倍系数和衰减率的确切值
\end{itemize}
攻击者即使完全了解算法，也无法精确计算裁定阈值，从而丧失了精确投毒操纵的能力。
\end{definition}

\begin{definition}[扰动层]
在每次攻防演练或重大安全事件复盘后，对秘密参数施加微小随机扰动：
\[
p_{\text{new}} = p_{\text{old}} \times (1 + \epsilon), \quad \epsilon \in [-\delta, \delta],
\]
其中扰动幅度 \(\delta\) 本身也是秘密。该机制等同于密码学中的密钥轮换，使攻击者的历史观测信息强制过期。
\end{definition}

\subsection{攻击者处境的转变}

在柯克霍夫原则下，攻击者面对的是一个透明但不可精确预测的系统。他想要用假告警操纵A值，但不知道下限在哪里，制造再多假告警也无法保证压到触发忽略的阈值。

\begin{center}
\begin{tabular}{c|c}
\toprule
\textbf{攻击者知道的} & \textbf{攻击者不知道的} \\
\midrule
优化公式的全部逻辑 & \(A_i^{\min}\) 的确切值 \\
反馈更新的类型 & 衰减触发的确切阈值 \\
每日配额分配的规则 & 外部威胁情报源的权重 \\
历史告警流和裁定结果 & 当前周期的扰动参数 \\
\bottomrule
\end{tabular}
\end{center}

\section{2.2 形式化安全定义：$\varepsilon$-参数不可区分性}

\subsection{从IND-CPA到参数不可区分性}

密码学中的IND-CPA定义了选择明文攻击下的不可区分性。我们将此范式迁移至告警分级领域，定义“参数不可区分性”。

\begin{definition}[\(\varepsilon\)-参数不可区分性]
防御方维护两个秘密参数集合 \(\theta_0\) 和 \(\theta_1\)。每次攻击方发起试探告警时，防御方随机选择 \(b \in \{0,1\}\)，基于 \(\theta_b\) 裁定告警并返回结果。攻击方观察结果，猜测 \(b\)。

若对任何计算资源有限的攻击方，
\[
\Pr[\text{攻击方猜对 } b] \le \frac{1}{2} + \varepsilon,
\]
且 \(\varepsilon\) 是可通过调整参数控制的小量，则称防御方案满足 \(\varepsilon\)-参数不可区分性。
\end{definition}

这个定义的价值在于：我们不要求攻击方完全猜不出参数，只要求他无法以高置信度区分两种不同的配置。这直接建立在实际的攻击方探测能力之上。

\subsection{第一层保障：信息论下界（统计防御）}

攻击方的核心操作是估计防御方响应概率 \(P_{\text{obs}}\)。由Chernoff-Hoeffding不等式，攻击方要将估计误差控制在 \(\Delta\) 以内，所需探测次数 \(m\) 的下界为：
\[
m \ge \frac{1}{2\Delta^2} \ln\frac{2}{\delta},
\]
其中 \(\delta\) 为置信度。若防御方将 \(\theta_0\) 和 \(\theta_1\) 的差异 \(\Delta\) 控制在0.01，攻击方需进行约38000次探测。这远超其在单周期内的资源预算（\(M_{\max}\)）。

\textbf{跨周期累积无效性}：由于参数在每个周期末被独立扰动，攻击方在不同周期收集的探测数据不具备相关性。他在新周期必须从头开始探测，之前的样本全部作废。只要 \(M_{\text{period}} < m^*\)，攻击方就永远无法完成精确推断。

\subsection{第二层保障：欠定推断（代数防御）}

蓝队可以主动注入噪声，以概率 \(q_{01}\) 将忽略伪装成分析，以概率 \(q_{10}\) 将分析伪装成忽略。攻击方观测到的响应概率为：
\[
P_{\text{obs}}(y=1) = r_i (1 - q_{10}) + (1 - r_i) q_{01}.
\]
攻击方面对三个未知量（\(r_i, q_{01}, q_{10}\)），但只有一个独立方程。这是一个欠定问题，在代数上不可能唯一确定参数值。

两层保障互补：统计上难以获得足够样本，代数上即使获得无限样本也无法求得唯一解。两者叠加，构成了 \(\varepsilon\)-参数不可区分性的完整数学基础。

\section{2.3 探测与反探测博弈}

\subsection{从被动采样到主动博弈}

传统的安全评估将攻击方的探测视为从自然中采样。但在对抗性环境中，蓝队是智能的、会主动扭曲样本的对手。探测与反探测本质上是一个非完全信息动态博弈。

\begin{definition}[探测与反探测博弈]
\textbf{攻击方}：选择探测序列，目标是最大化对参数 \(\mathbf{r}\) 的推断准确性。
\textbf{防御方}：在每个被探测的告警类型上，选择伪造策略（\(q_{01}, q_{10}\)），最大化攻击方的推断误差。蓝队受到反探测预算 \(C_{\text{counter}}\) 的约束。
\end{definition}

\subsection{蓝队最优反探测策略}

蓝队的目标是使攻击方的观测样本信噪比最小化。在反探测预算有限的情况下，蓝队的资源分配遵循以下原则：
\begin{itemize}
    \item \textbf{模糊地带强投入}：对 \(r_i \approx 0.5\) 的类型，攻击方最不确定。蓝队投入少量反探测资源就能极大增加其推断难度。
    \item \textbf{极端伪装}：对 \(r_i\) 接近0或1的类型，蓝队只需投入极少资源制造假阳性或假阴性，就能颠覆攻击方的判断。
\end{itemize}

\subsection{蜜罐式反探测与错觉塑造}

攻击方会优先探测那些“历史上最少被攻击”的冷门类型。蓝队可以反向利用这一点：
\textbf{蜜罐式反探测}：在攻击方可能探测的冷门方向上，故意制造少量“被分析”的伪造历史记录，让攻击方误以为这些地方重兵把守。这能将攻击方的火力引向蓝队真正布防的区域，实现了对攻击方信念的主动塑造。

\section{诚实标注}

\begin{center}
\begin{xltabular}{\textwidth}{c|X|c}
\toprule
\textbf{命题} & \textbf{状态} & \textbf{层级} \\
\midrule
柯克霍夫三层架构 & 框架完整，逻辑自洽，与AES设计哲学同构 & II（跨领域适用性待验证） \\
参数不可区分性定义 & 形式化定义已建立，术语已明晰 & II（安全参数的量化待社区共识） \\
信息论下界（Chernoff-Hoeffding） & 基于标准的概率不等式，推导严格 & I（不依赖未证明假设） \\
欠定推断的不可能性 & 方程组系数矩阵的秩分析，结论严格 & I（严格数学结论） \\
最优反探测策略 & 定性结论已给出，符合直觉和博弈论原理 & II（精确解待实验数据标定） \\
蜜罐式反探测 & 概念已形式化，战术启发性强 & III（实战效果待红蓝对抗检验） \\
\bottomrule
\end{xltabular}
\end{center}

\vspace{5mm}
\noindent\textbf{文档信息}：
\begin{itemize}
    \item 作者：李龙胜
    \item 版本：v1.0（四卷体系重构第2卷）
    \item 许可：CC BY 4.0
    \item 前置依赖：第0卷、第1卷
\end{itemize}

\end{document}